\doublespacing % Do not change - required

\chapter{User and Deployment Guide}
\label{app:userguide}

\thesisspacing % Do not change - required

This appendix is a concise guide to reproducing the work and deploying the translator. The complete code, configurations and trained models are in the project repository.

\section{Repository}
\label{app:repo}

The complete code and configurations are held in a private Git repository, with the trained checkpoints stored alongside it, and access to both is available to the supervisor and examiners on request. The layout follows the experiment structure used throughout this report, with training and distillation configurations and evaluation outputs under \path{experiments/}, export and benchmark scripts under \path{scripts/}, and per-experiment run scripts in the repository root. The selected deployment student is \path{experiments/student_l1b_c_featkd}.



\section{Environment Setup}
\label{app:setup}

Two environments are used. Training, export and evaluation run on a Linux workstation with an NVIDIA GPU, using PyTorch with the frozen foundation models (DINOv2, CLIP, SegFormer, DepthAnything) pulled from their public hubs, plus ONNX Runtime and Google's litert-torch for the CPU export paths. Compilation to the Hailo Executable Format uses the Hailo Dataflow Compiler 3.33.1 on Python 3.10, which is distributed only inside the vendor Docker image and is run from there (Section~\ref{sec:impl_hailo}). The Raspberry~Pi~5 runs Raspberry Pi OS with HailoRT 4.23.0 and the \texttt{hailortcli} tool for the Hailo-8L HAT, and picamera2 for the capture rig.


\section{Reproducing the Results}
\label{app:reproduce}

The teacher is trained first, then the width-16 student is distilled from it with the objective of Section~\ref{sec:impl_student}, and the student is exported to LiteRT, ONNX and HEF. All quality and downstream numbers in Chapter~\ref{ch:results} are regenerated by the evaluation harness, which compares raw NIR, native RGB and translated RGB through the frozen model panel and writes per-image metrics, so the tables and confidence intervals can be rebuilt without rerunning every model. The Raspberry~Pi~5 CPU benchmark is reproduced with:
\begin{lstlisting}[language=bash]
python3 scripts/benchmark_pi_cpu.py --models_dir experiments/student_l1b_c_featkd
\end{lstlisting}


\section{Deploying to the Raspberry Pi 5}
\label{app:deploy}

There are two deployment paths. For the CPU path, copy the 2.3\,MB mixed-INT8 LiteRT artefact to the Pi and run it through the LiteRT XNNPACK runtime, no accelerator is needed. For the NPU path, compile the FP32 ONNX student to a HEF with the Dataflow Compiler as in Section~\ref{sec:impl_hailo}, copy the resulting \path{student_small_nir_to_rgb.hef} to the Pi, and run it on the Hailo-8L with \texttt{hailortcli}:
\begin{lstlisting}[language=bash]
hailortcli run student_small_nir_to_rgb.hef --measure-latency
\end{lstlisting}
The live demonstration of Section~\ref{sec:res_live_demo} is served by a small web application that runs as a systemd service on the Pi, defaults to the Hailo HEF backend, and is reachable in a browser on port 8080. The translator backend can be switched at runtime between the Hailo HEFs and the CPU LiteRT profiles. One practical requirement: the Hailo HAT draws enough additional current that an adequately rated power supply is needed to avoid brown-outs under sustained load.


